                \begin{tikzpicture}[xscale=1.0, yscale=1.0]
                    \useasboundingbox (-0.35, -2.05) rectangle (12.6, 3.35);

                    \tikzset{jeton/.style={draw=inria-2024-bleu-canard, line width=0.7pt, rounded corners=1.5pt,
                                           fill=inria-2024-bleu-canard!12, inner sep=2pt, font=\scriptsize,
                                           minimum height=4.2mm}}
                    \tikzset{modele/.style={draw=gris_fonce_inria, line width=1.0pt, rounded corners=3pt,
                                            fill=gray!8, inner sep=4pt, font=\small\bfseries, minimum width=1.55cm,
                                            minimum height=6.5mm}}
                    \tikzset{modeleouvert/.style={modele, draw=inria-rouge, fill=inria-rouge!6, text=inria-rouge}}
                    \tikzset{fl/.style={-{Latex[length=2.2mm]}, line width=0.9pt, color=gray!55}}
                    \tikzset{dom/.style={font=\small, anchor=east, text=gris_fonce_inria}}

                    % ---------------- ligne 1 : le langage ----------------
                    \node[dom] at (1.35, 2.75) {\textbf{langage}};
                    \node[jeton, anchor=west] (t1) at (1.75, 2.75) {Les};
                    \node[jeton, anchor=west] (t2) at ($(t1.east)+(0.09,0)$) {graph};
                    \node[jeton, anchor=west] (t3) at ($(t2.east)+(0.09,0)$) {es};
                    \node[jeton, anchor=west] (t4) at ($(t3.east)+(0.09,0)$) {sont};
                    \node[jeton, anchor=west] (t5) at ($(t4.east)+(0.09,0)$) {partout};
                    \node[font=\scriptsize, text=inria-2024-bleu-canard, anchor=west] at ($(t5.east)+(0.25,0)$) {sous-mots};
                    \node[modele, anchor=west] (m1) at (9.85, 2.75) {GPT};
                    \draw[fl] (8.55, 2.75) -- (m1.west);

                    % ---------------- ligne 2 : l'image ----------------
                    \node[dom] at (1.35, 1.45) {\textbf{image}};
                    \begin{scope}[shift={(1.80,0.80)}]
                        \foreach \i in {0,...,4} {
                            \foreach \j in {0,...,2} {
                                \draw[inria-2024-bleu-canard, line width=0.6pt, fill=inria-2024-bleu-canard!12]
                                    (\i*0.44, \j*0.44) rectangle ++(0.40,0.40);
                            }
                        }
                    \end{scope}
                    \node[font=\scriptsize, text=inria-2024-bleu-canard, anchor=west] at (4.25, 1.45) {imagettes};
                    \node[modele, anchor=west] (m2) at (9.85, 1.45) {SAM~2};
                    \draw[fl] (8.55, 1.45) -- (m2.west);

                    % ---------------- ligne 3 : la 3D ----------------
                    \node[dom] at (1.35, 0.05) {\textbf{monde 3D}};
                    \begin{scope}[shift={(1.80,-0.55)}]
                        \foreach \x/\y in {0.05/0.10, 0.22/0.52, 0.38/0.86, 0.55/0.30, 0.74/0.68, 0.92/1.02,
                                           1.10/0.16, 1.28/0.60, 1.44/0.94, 1.62/0.38, 1.80/0.76, 1.98/0.22,
                                           0.15/0.72, 0.66/0.98, 1.20/0.92, 1.72/1.04, 0.44/0.44, 0.98/0.48} {
                            \fill[inria-2024-bleu-canard] (\x, \y) circle (1.6pt);
                        }
                    \end{scope}
                    \node[font=\scriptsize, text=inria-2024-bleu-canard, anchor=west] at (4.25, 0.05) {points du capteur};
                    \node[modeleouvert, anchor=west] (m3) at (9.85, 0.05) {?};
                    \node[font=\scriptsize, text=gris_fonce_inria, anchor=north, align=center] at (10.62, -0.30)
                        {Sonata, Concerto,\\Utonia\ldots};
                    \draw[fl, color=inria-rouge!55, densely dashed] (8.55, 0.05) -- (m3.west);

                    % ---------------- séparateur ----------------
                    \draw[gray!35, line width=0.6pt] (0.55, 0.62) -- (11.9, 0.62);

                    % ---------------- accolade et question ----------------
                    \node[font=\scriptsize, align=center, text=inria-rouge, anchor=north] at (5.4, -1.05)
                        {en 3D, l'unité de calcul reste celle de l'\textbf{échantillonnage} :\\
                         point du capteur ou case de grille};
                \end{tikzpicture}
